problem:  
Let \(X\) be a compact complex threefold satisfying the Hull–Strominger system, endowed with a closed 3‑form \(H\) whose integral cohomology class \([H] \in H^3(X,\mathbb{Z})\) may have torsion.  Denote by \(\mathcal{G}_{\widehat{H}}\) a **bundle gerbe with connection and curving** whose Dixmier–Douady class in \(H^3(X,\mathbb{Z})\) represents the differential cohomology class \(\widehat{H} \in \widehat{H}^3(X,\mathbb{Z})\).  

It is known that smooth exact Courant algebroids correspond to fractional Ševera classes in \(H^3(X,\mathbb{R})\).  For a gerbe \(\mathcal{G}_{\widehat{H}}\), one can construct a **gerbe‑twisted Courant algebroid** \(E_{\mathcal{G}_{\widehat{H}}}\) by performing local \(B\)-field gluings on \(TX \oplus T^*X\) using the gerbe’s connective data; this refines the ordinary \(H\)-twisted Courant algebroid by the integral and torsion information of \(\mathcal{G}_{\widehat{H}}\).

> **Problem:**  
> For a fixed gerbe \(\mathcal{G}_{\widehat{H}}\) on a compact complex manifold \(X\):  
> 1. Construct a **deformation complex** governing infinitesimal deformations of generalized complex structures compatible with \(\mathcal{G}_{\widehat{H}}\)’s connective structure; this must reduce to the usual deformation complex when the gerbe is topologically trivial.  
> 2. Investigate whether the torsion part of the Dixmier–Douady class of \(\mathcal{G}_{\widehat{H}}\) can contribute nontrivially to the local analytic moduli of solutions of the Strominger system realized within \(E_{\mathcal{G}_{\widehat{H}}}\).  
> 3.  Concretely, does there exist a pair of gerbes \(\mathcal{G}_1,\mathcal{G}_2\) with distinct torsion Dixmier–Douady classes whose associated generalized complex moduli spaces \(\mathcal{M}_{\mathcal{G}_i}\) have different analytic germ structure near a common underlying \(SU(3)\)-structure?  

Why is it a "good" problem:  
This problem is carefully defined within established geometric frameworks—bundle gerbes with connection and their Courant extensions—avoiding ill‑defined objects while keeping torsion data explicit.  It targets a deep and still unresolved issue: whether **discrete topological torsion**, normally invisible to smooth deformation theory, can manifest analytically once one couples generalized geometry to gerbe connections.  It bridges **differential cohomology**, **gerbe theory**, and **non‑Kähler deformation problems** in heterotic string geometry, potentially revealing new analytic and topological interactions.  The statement is simple yet potent: it asks whether torsion in a gerbe’s class can alter analytic moduli—a question at the frontier of higher geometry and mathematical physics, satisfying all hallmarks of a “good’’ research problem.